\documentclass[11pt]{letter}
% \usepackage[margin=1in]{geometry}
\usepackage[left=1cm,right=1cm,top=1.8cm,bottom=1.5cm]{geometry}
\usepackage[T1]{fontenc}
\usepackage{times}
\usepackage{url}

\signature{Julien Soulé\\On behalf of all authors}
\address{Julien Soulé\\SEDAN Research Group, SnT\\University of Luxembourg}

\begin{document}

\begin{letter}{Editor-in-Chief\\IEEE Transactions on Cloud Computing}

  \opening{Dear Editor-in-Chief,}

  We submit the manuscript entitled ``KARMA: An Automated Multi-Agent Framework for Resilient Kubernetes Autoscaling'' for consideration as a Regular Paper in IEEE Transactions on Cloud Computing. We request traditional publication rather than Open Access. The manuscript is not under review elsewhere, and all authors have approved this submission.

  This submission is an extended journal version of the IEEE CLOUD 2025 paper ``Streamlining Resilient Kubernetes Autoscaling with Multi-Agent Systems via an Automated Online Design Framework'' by Soulé, Jamont, Occello, Traonouez, and Théron. The manuscript explicitly cites the earlier paper and identifies it as a preliminary version. To the best of our knowledge, this submission does not violate any copyright commitments associated with the conference publication.

  The TCC submission contains substantial additional technical material beyond the preliminary CLOUD version. The main extensions are:

  \begin{itemize}
    \item a more explicit organizationally guided MARL formulation for resilient Kubernetes autoscaling, including role-guided action constraints and mission-guided reward shaping;
    \item a revised KARMA framework description covering digital-twin construction, MAPPO training, behavioral analysis, and guarded transfer to Kubernetes;
    \item expanded experimental methodology with reproducibility controls, baseline alignment, and six research questions;
    \item additional evaluations covering ablations, disturbance-intensity robustness, sim-to-real gap analysis, behavioral separability, and computational cost;
    \item a broader discussion of deployment conditions, scalability, role engineering validity, security validity, and reproducibility limitations;
    \item an updated related-work section and bibliography, including recent work on Kubernetes/RL autoscaling, cloud workload autoscaling, and MARL-based scaling.
  \end{itemize}

  Together, these additions exceed the 35\% additional technical material requirement for extended versions. The journal manuscript reframes the contribution around organizational guidance as a mechanism for making learned autoscaling policies more inspectable, bounded, and auditable for cloud operators.

  Supplementary material is submitted separately from the main manuscript. It summarizes the artifact contents expected for reproduction, including trace schemas, configurations, hyperparameter settings, run logs, scripts, and secondary figures referenced from the paper.

  \closing{Sincerely,}

\end{letter}

\end{document}
